\documentclass[sigconf,article]{acmart}

% for NTCIR proceedings
\settopmatter{printacmref=false} % Removes citation information below abstract
\renewcommand\footnotetextcopyrightpermission[1]{} % removes footnote with conference information in first column.
\pagenumbering{gobble} % removes the page numbers for publication

\usepackage{makecell}

\begin{document}

\title{Overview of the NTCIR-19 AgenticInstruction Task}

\author{Hideo Joho}
\affiliation{%
  \institution{University of Tsukuba}
  \city{Tsukuba}
  \country{Japan}
}
\email{hideo@slis.tsukuba.ac.jp}

\begin{abstract}
This paper describes the overview of the NTCIR-19 AgenticInstruction task.
This pilot task aims to identify effective and robust instructions for guiding
Large Language Models (LLMs) in search agent roles. Participants were invited
to submit instructions and generation methods covering four key search stages:
query formulation, document selection (clicks), relevance judgement, and query
reformulation. AgenticInstruction-1 centred on a high-recall ad hoc retrieval
scenario in both Japanese and English. Submitted instructions were evaluated
in an iterative search setting with performance measured at the session level:
the number of unique relevant documents identified (recall) against the number
of tokens consumed by the LLM.
% TODO after run submissions (2026-08-10):
The task attracted \textbf{[N]} teams which submitted a total of
\textbf{[M]} runs. In this paper we describe the task design, the evaluation
framework, datasets, organiser baselines, and the results of submitted runs
with meta analyses.
\end{abstract}

\keywords{agentic search, instruction design, large language models,
evaluation, high-recall retrieval, test collections}

\maketitle

\section{Introduction}

Large Language Models (LLMs) are increasingly deployed as autonomous search
agents that formulate queries, examine retrieved documents, judge their
relevance, and reformulate queries over multiple iterations. In such systems,
the natural-language \emph{instructions} that direct each stage of the search
process have become a decisive factor in overall performance. However, no
systematic comparison of instruction designs exists for search agents: prior
work has largely examined single search stages in isolation, whereas real
search is multi-stage and iterative. The AgenticInstruction task evaluates
instructions in full search sessions, making the \emph{instruction} --- rather
than the retrieval system --- the object of submission and comparison.

The task addresses the following research questions:

\begin{itemize}
  \item \textbf{RQ1a}: What are the best practices for instruction design in
  agentic search?
  \item \textbf{RQ1b}: What characterises effective and robust instructions
  for query formulation, document selection, relevance judgement, and query
  reformulation?
  \item \textbf{RQ1c}: How does instruction performance evolve over multiple
  iterations in the search process?
  \item \textbf{RQ2a}: What is the relationship between instruction design
  features and ranking models?
  \item \textbf{RQ2b}: What is the relationship between instruction design
  features and the choice of LLMs?
  \item \textbf{RQ2c}: To what extent do search topics influence instruction
  performance?
\end{itemize}

As a pilot task, AgenticInstruction-1 did not take up RQ2b at the official-run
level: a single LLM was fixed for all submitted runs to maximise comparability
across participants (Section~\ref{sec:framework}), and the influence of LLM
choice was delegated to additional analyses in participant papers.

The remainder of this paper is organised as follows.
Section~\ref{sec:task} presents the task design, including changes made since
the initial design presented at the NTCIR-19 kick-off event.
Section~\ref{sec:framework} describes the evaluation framework.
Section~\ref{sec:datasets} introduces the datasets.
Section~\ref{sec:baselines} reports the organiser baselines.
Section~\ref{sec:results} presents the participation and results of submitted
runs, and Section~\ref{sec:meta} provides meta analyses.
Section~\ref{sec:conclusions} concludes the paper.

\section{Task Design}
\label{sec:task}

\subsection{Search Task}

AgenticInstruction-1 employs a \emph{high-recall ad hoc retrieval} scenario:
find as many different relevant documents as possible for a given search topic
from a given document collection using a provided search tool. An LLM-based
search agent performs a session of up to five \emph{iterations}, where each
iteration consists of four stages:

\begin{enumerate}
  \item \textbf{Query formulation} (first iteration) or \textbf{query
  reformulation} (subsequent iterations): the LLM produces a query string;
  \item \textbf{Ranking}: the search tool returns a result page (SERP) of ten
  documents;
  \item \textbf{Document selection (click)}: the LLM selects the documents
  whose full text should be examined, or none;
  \item \textbf{Relevance judgement}: the LLM judges each selected document as
  relevant or not relevant based on its full text.
\end{enumerate}

The goal of the task is to develop instruction methods for an LLM to
\emph{maximise the total number of unique relevant documents identified
(recall)} within five iterations, \emph{while minimising the total number of
tokens} sent to and produced by the LLM. A document counts as relevant for
evaluation only if it satisfies all three of the following conditions:
(1) it was clicked (selected) from a SERP by the LLM;
(2) it was judged as relevant by the LLM; and
(3) it is labelled as relevant by the official qrels.

\subsection{Instructions as Submissions}

Participants design the instruction text given to the LLM at each stage.
Figure~\ref{fig:template} shows the instruction template for the query
formulation stage: the participant-supplied instruction is combined with fixed
descriptions of the task, corpus, and search tool, and the topic description.
Customisable parameters include the per-stage instructions, the system prompt,
the topic fields presented to the model (title, description, narrative), and
the search tool description. The task and corpus descriptions are fixed for
submitted runs.

\begin{figure}
\begin{verbatim}
Instruction: {participant instruction}
============================
Task Description: {task description}
Corpus Description: {corpus description}
Search Tool Description: {tool description}
Topic Description:
- Title: {title}
- Description: {description}
- Narrative: {narrative}
\end{verbatim}
\caption{Instruction template for the query formulation stage.}
\label{fig:template}
\end{figure}

\subsection{Run Types}

Following NTCIR conventions, two run types were accepted:

\begin{itemize}
  \item \textbf{Automatic run}: any training or human intervention is allowed
  on the train/dev datasets, but no human intervention is allowed on the test
  datasets; the instruction method must be frozen before the test datasets are
  used.
  \item \textbf{Manual run}: human intervention is allowed on the test
  datasets as well.
\end{itemize}

Each group could submit up to three runs, where one run consists of session
log files for at least two configurations (BM25 and Sparse Encoder) in the
participant's choice of language(s).

\subsection{Changes from the Initial Design}
\label{sec:changes}

Several aspects of the task were revised between the kick-off event
(September 2025) and the formal run period, primarily to reduce technical
barriers and to maximise comparability across participants:

\begin{itemize}
  \item \textbf{Distribution}: a planned Docker image was cancelled due to
  technical and licensing issues; instead, the organisers hosted a shared
  OpenSearch service with per-group credentials, and provided step-by-step
  environment setup instructions.
  \item \textbf{LLMs}: the initial design listed multiple open-weight models
  (gpt-oss, Llama, Qwen); the official runs were fixed to a single model,
  \texttt{gpt-oss-120b}, accessed through Amazon Bedrock. Participants were
  welcome to test other LLMs in additional analyses.
  \item \textbf{IR models}: DPR was dropped from the official configuration,
  leaving BM25 and a sparse encoder (SPLADE-style) as the two official
  ranking models.
  \item \textbf{LLM access}: the task initially adopted a free-tier LLM API
  service (Groq); as free-tier token limits proved insufficient for full
  sessions and new paid-tier registrations were closed, the official access
  path was switched to Amazon Bedrock.
\end{itemize}

\section{Evaluation Framework}
\label{sec:framework}

\subsection{geniie-lab}

All experiments are conducted with \texttt{geniie-lab}\footnote{%
\url{https://github.com/geniie-lab/geniie-lab}}~\cite{joho2025sigirap}, an
open-source controlled experimentation testbed for model search behaviour,
which operationalises the instruction--response perspective on LLMs in
information retrieval tasks proposed by Joho and Jose~\cite{joho2025sigir}. The testbed
instructs an LLM to perform search actions (querying, clicking, judging
relevance, reformulating) across configurable, modular stages, records every
stage output as structured JSONL logs, and computes session-level statistics.
Participants run their instruction methods locally with \texttt{geniie-lab}
and submit the generated log files; the submitted logs are the run.

\subsection{Fixed Components}

For submitted runs, the following components were fixed:

\begin{itemize}
  \item \textbf{LLM}: \texttt{gpt-oss-120b} (open-weight, 117B-parameter
  mixture-of-experts model) accessed via Amazon Bedrock, with greedy
  decoding settings (temperature 0.0, top-$p$ 1.0);
  \item \textbf{Search tools}: OpenSearch with (a) BM25 ranking (with
  language-appropriate analysis, e.g.\ Kuromoji for Japanese) and (b) a
  multilingual learned sparse encoder\footnote{%
  \path{opensearch-neural-sparse-encoding-multilingual-v1}} over
  passage-chunked documents, hosted by the organisers;
  \item \textbf{SERP size}: ten results per query;
  \item \textbf{Session plan}: five iterations of the four-stage sequence.
\end{itemize}

\subsection{Measures}

The primary measures are computed per topic from the session logs:

\begin{itemize}
  \item \textbf{Unique relevant documents}: the number of distinct documents
  satisfying the triple condition (clicked, LLM-judged relevant, and
  qrels-relevant), accumulated across iterations;
  \item \textbf{Recall}: unique relevant documents divided by all relevant
  documents in the official qrels for the topic;
  \item \textbf{Total tokens}: the sum of tokens sent to and produced by the
  LLM across all stages of the session;
  \item \textbf{Per-iteration progression}: the cumulative unique relevant
  documents at the end of each iteration, showing where improvements level
  off.
\end{itemize}

Runs are compared primarily on a \emph{recall/token} basis: for a given token
budget, higher recall is better; for a given recall, fewer tokens are better.

\subsection{Reproducibility Considerations}

Although all official runs use greedy decoding (temperature~0), large hosted
models are not exactly reproducible: serving-side effects (batching,
mixture-of-experts routing, kernel non-determinism) cause small run-to-run
variations that compound over a 20-stage session. In repeated organiser
baseline runs, 50-topic mean recall varied by roughly $\pm 0.015$ between
identical configurations. This inter-run variance constitutes a noise floor
against which differences between instruction methods should be interpreted.

\section{Datasets}
\label{sec:datasets}

Table~\ref{tab:datasets} summarises the datasets. English uses newswire
collections from TREC Robust; Japanese uses academic abstract collections
from NTCIR-1 and NTCIR-2. Train/dev datasets were available throughout the
task period; participants were expected to obtain test datasets only when
they decided to submit a run, and all development had to be carried out on
the train/dev datasets. For the Japanese qrels, both grades ``relevant'' and
``partially relevant'' count as relevant (qrels label~$\geq 1$).

\begin{table}
\caption{Datasets of the AgenticInstruction task.}
\label{tab:datasets}
\begin{tabular}{@{}lllr@{}}
\toprule
Language & Role & Dataset & Topics \\
\midrule
English & Train/Dev & TREC Robust 2004 (folds 1--2) & 100 \\
English & Test & TREC Robust 2005 & 50 \\
Japanese & Train/Dev & NTCIR-1 AdHoc & 83 \\
Japanese & Test & NTCIR-2 AdHoc & 50 \\
\bottomrule
\end{tabular}
\end{table}

Documents were indexed by the organisers into per-collection OpenSearch
indexes: a BM25 index with language-appropriate analyzers, and a sparse
index in which documents are chunked into passages server-side and each
passage is encoded into a learned sparse representation; at query time, a
document is scored by its best-matching passage. Queries are encoded on the
participant side with the same multilingual sparse encoder, so no
model-serving permissions are required of participants.

\section{Organiser Baselines}
\label{sec:baselines}

To provide participants with a reference point, the organisers produced
baselines for all six train/dev configurations using the default
instructions distributed with the task (e.g.\ ``Review the provided
descriptions of task, corpus, tool and search topic. Then, formulate a search
query.''), with no optimisation. Table~\ref{tab:baselines} summarises the
results; the full per-topic tables are published on the task
website.\footnote{\url{https://geniie-lab.github.io/ntcir/}}

\begin{table}
\caption{Organiser baselines (default instructions, five iterations,
\texttt{gpt-oss-120b} via Amazon Bedrock).}
\label{tab:baselines}
\begin{tabular}{@{}llrrr@{}}
\toprule
Dataset & IR model & \makecell[r]{Mean uniq.\\rel.\ docs} & \makecell[r]{Mean\\recall} & \makecell[r]{Total\\tokens} \\
\midrule
Robust 2004 fold1 & BM25 & 4.44 & 0.127 & 34.0M \\
Robust 2004 fold1 & Sparse & 4.98 & 0.154 & 44.2M \\
Robust 2004 fold2 & BM25 & 4.52 & 0.134 & 30.2M \\
Robust 2004 fold2 & Sparse & 3.86 & 0.118 & 30.2M \\
NTCIR-1 AdHoc & BM25 & 6.45 & 0.239 & 51.0M \\
NTCIR-1 AdHoc & Sparse & 4.93 & 0.188 & 40.8M \\
\bottomrule
\end{tabular}
\end{table}

Three observations characterise the default-instruction behaviour. First,
\emph{improvements concentrate in early iterations}: roughly 75--80\% of
final recall is obtained in the first iteration, and gains largely level off
by the third iteration, while later iterations are the most token-expensive
because the conversation history is at its largest. Second, in English the
two IR models perform comparably under default instructions, whereas in
Japanese BM25 clearly outperforms the multilingual sparse encoder (mean
recall 0.239 vs.\ 0.188). Third, session cost is dominated by clicking
behaviour: the default instructions click generously (often most of the
SERP), so a full session consumes roughly 0.6--0.9M tokens per topic.
Together these define the challenge posed to participants: make later
iterations productive, or stop early and save tokens.

\section{Participation and Results}
\label{sec:results}

% TODO after run submissions (2026-08-10) and evaluation (2026-09-01):
% - Table: participating teams (Team ID, name, country, language(s), #runs)
% - Table: run results per configuration (mean recall, total tokens),
%   ordered by recall; include organiser baseline rows for reference
% - Figure: recall/token scatter per configuration (one point per run;
%   baseline highlighted)
% - Figure or table: per-iteration progression of the best runs vs baseline

\textbf{[To be completed after run submissions close on 10 August 2026.]}

\section{Meta Analyses}
\label{sec:meta}

% TODO after participant papers (draft due 2026-10-01):
% - RQ1a/RQ1b: taxonomy of instruction strategies across teams per stage
% - RQ1c: iteration progression comparison across runs
% - RQ2a: instruction x IR model interaction (paired runs)
% - RQ2c: per-topic difficulty analysis (topic variance, hard topics)
% - Token/recall efficiency analysis

\textbf{[To be completed after run submissions and participant paper
drafts.]}

\section{Conclusions and Future Directions}
\label{sec:conclusions}

This paper presented an overview of the NTCIR-19 AgenticInstruction task, a
pilot task for identifying effective and robust instructions for LLM-based
search agents, evaluated at the session level in a high-recall ad hoc
scenario across English and Japanese.
% TODO: summarise participation and principal findings once available.

Several directions are candidates for a future cycle: graded relevance in
the session evaluation; additional languages and genres; letting the agent
plan its own session (agentic plans) rather than following a fixed
iteration structure; instruction robustness protocols across serving stacks;
and multi-LLM comparisons (RQ2b), which this pilot deliberately left to
participant analyses.

\begin{acks}
% TODO: participants, dataset providers (NII for NTCIR collections, NIST/LDC
% for TREC collections), infrastructure support.
\end{acks}

\bibliographystyle{ACM-Reference-Format}
\bibliography{overview}

\end{document}
